\section{A bounded induction difference in the cubic codegree-three branch}
\label{sec:cubic-codegree-three-difference}
Consider a self-dual pseudo-unitary rank-four category whose rational Grothendieck
algebra is $\mathbb Q\times K$, with $K$ cyclic cubic, and whose rational
character has formal codegree three. Normalize this character to
$\rho=(1,0,1,-1)$. In the resulting basis $(1,X,Y,Z)$ the tensor tuple
has the form
\[
 (a,b,b+1,d,d,d,g,g,g+1,g+1),\qquad
 ad+(2b+1)g-b^2-2d^2+1=0.
\]
Let $x,y,z$ be the Frobenius--Perron dimensions. Character orthogonality
gives $z=y+1$. The cubic character hypothesis gives the unique rational
relation $(1,0,1,-1)$ among $1,x,y,z$.

Write $D=\operatorname{FPdim}\mathcal C$ and
$I(1)=1\oplus B\oplus A\oplus A'$, where $B$ corresponds to $\rho$.
Its dimension is $D/3$. Comparison of its restriction dimension in the
cubic field forces
\[
 F(B)=(1,p,q,q+1),\qquad
 p=\frac{a+2d}{3},\quad q=\frac{b+2g}{3},\qquad
 p,q\in\mathbb Z_{\geq0}.
\]
In particular the indicated congruences modulo three are necessary.

\begin{proposition}\label{prop:cubic-bounded-difference}
There is an integer $k\in\{-1,0,1\}$ such that, on the two restrictions
$F(A),F(A')$, the multiplicity differences $[Y]-[Z]$ are $k,-k$.
For each remaining central simple $W$, put
\[
 h_W=[F(W):Z]-[F(W):Y].
\]
Then
\[
 \sum_{W\notin\{1,B,A,A'\}}h_W^2=6-2k^2.
\]
Thus at most six remaining central simples can have different
$Y$- and $Z$-multiplicities, independently of the fusion coefficients.
If $k=0$, the multiset of nonzero $|h_W|$ is either six ones or
$(2,1,1)$. If $|k|=1$, it is either four ones or $(2)$.
\end{proposition}
\begin{proof}
Let $M$ be the full restriction matrix of central induction, and
$Q=\sum_U N_U^2=MM^{\mathsf T}$. The first column of $Q$ is
\[
 (4,3p,3q+1,3q+2)^{\mathsf T}.
\]
Since precisely the four constituents of $I(1)$ have unit coefficient
one in their restrictions, the restriction columns of $A,A'$ sum to
$(2,2p,2q+1,2q+1)^{\mathsf T}$. Their $Y-Z$ differences therefore
sum to zero; denote them by $k,-k$.

Put $v=(1,0,1,-1)^{\mathsf T}$ and
$w=(0,0,-1,1)^{\mathsf T}=e_0-v$. Since $\rho$ has codegree three,
$Qv=3v$ and $v^{\mathsf T}v=3$. As $Q_{00}=4$, we obtain
\[
 w^{\mathsf T}Qw=4-2\cdot3+3\cdot3=7.
\]
The four known columns contribute $0,1,k^2,k^2$ to this Gram expression.
Subtracting proves the displayed identity. Its nonnegative integer left
side forces $k^2\leq3$, hence $k\in\{-1,0,1\}$. The final assertions
are the elementary decompositions of four or six as sums of nonzero
integer squares.
\end{proof}

There are corresponding moment identities involving at most six twists.
Put $\alpha=\dim A$, $\beta=\dim A'$ and $a_W=\dim W>0$. Then
\begin{align*}
 \sum_W h_Wa_W&=\frac{2D}{3}+k(\alpha-\beta),\\
 \sum_W h_Wa_W\theta_W&=-\frac D3+k(\alpha-\beta),\\
 \sum_W h_Wa_W\theta_W^2
 &=D\bigl(\nu_2(Z)-\nu_2(Y)\bigr)-\frac D3+k(\alpha-\beta),
\end{align*}
where the sums range only over the remaining central simples.
Indeed $\dim I(Z)-\dim I(Y)=D$ since $z-y=1$, while the known
twist-one contribution is $D/3-k(\alpha-\beta)$. Subtraction from
the first and second indicator identities proves the other two formulas.
More generally the last right side is valid with $\nu_n$ and $\theta^n$
for every $n\geq1$.

\begin{corollary}
If $|k|=1$, exactly four remaining central simples have nonzero
$Y-Z$ multiplicity difference, and each difference has absolute value one.
\end{corollary}
\begin{proof}
It remains to exclude the single-term pattern $(2)$. In that case the
zeroth and first moments above are real numbers of the same absolute
value, since the twist has absolute value one. They differ by $D$, so
they must be negatives of each other. Therefore
\[
 k(\alpha-\beta)=-D/6.
\]
Since $\alpha+\beta=2D/3-1$, one of $\alpha,\beta$ equals
$D/4-1/2$. The two dimensions are $D/\sigma(D)$ for the two nontrivial
automorphisms $\sigma$ of the cyclic cubic field. Consequently one such
automorphism satisfies
\[
 \sigma(1/D)=\frac14-\frac1{2D}.
\]
Iterating this affine identity three times gives $1/D=1/6$, impossible
because $D$ has degree three. For completeness, $D$ cannot be rational:
otherwise the three cubic characters would all have the same formal
codegree $D$, and the reciprocal-codegree identity would give
$1/3+3/D=1$, so $D=9/2$, contradicting algebraic integrality.
\end{proof}

\paragraph{What this reduction does not prove.}
The coefficients $h_W$ may be negative, and the positive dimensions
$a_W$ are not bounded by the number of nonzero terms. Thus these identities
are not sums of a bounded number of roots of unity with unit weights.
Neither the positive-measure inequalities nor a cyclotomic length bound
can be applied to this signed measure without an additional argument.
The reduction bounds the number of central objects that must be analyzed;
it does not yet bound the original fusion coefficients or classify this
cubic branch. No unrestricted categorification classification is inferred
from it.
